\documentclass[10pt,a4paper]{amsart}
\usepackage{amsmath,amssymb,amsthm,graphicx}
% Retained minimal preamble for the Stacks Project chapter flat.tex (More on Flatness).
% Verbatim excerpt of the chapter preamble: Xy-pic setup (native preamble lines 14-18) and the native
% theorem declarations with the mathematical macros (native preamble lines 154-217). Copied mechanically;
% nothing added, removed or redefined. The wrapper class and the amsmath/amssymb/amsthm/graphicx base
% packages are supplied by the assembler, so the class/externaldocument part of the native preamble is omitted.
% For commutative diagrams we use Xy-pic
\usepackage[all]{xy}

% We use 2cell for 2-commutative diagrams.
\xyoption{2cell}
% Theorem environments.
%
\theoremstyle{plain}
\newtheorem{theorem}[subsection]{Theorem}
\newtheorem{proposition}[subsection]{Proposition}
\newtheorem{lemma}[subsection]{Lemma}

\theoremstyle{definition}
\newtheorem{definition}[subsection]{Definition}
\newtheorem{example}[subsection]{Example}
\newtheorem{exercise}[subsection]{Exercise}
\newtheorem{situation}[subsection]{Situation}

\theoremstyle{remark}
\newtheorem{remark}[subsection]{Remark}
\newtheorem{remarks}[subsection]{Remarks}

\numberwithin{equation}{subsection}

% Macros
%
\def\lim{\mathop{\mathrm{lim}}\nolimits}
\def\colim{\mathop{\mathrm{colim}}\nolimits}
\def\Spec{\mathop{\mathrm{Spec}}}
\def\Hom{\mathop{\mathrm{Hom}}\nolimits}
\def\Ext{\mathop{\mathrm{Ext}}\nolimits}
\def\SheafHom{\mathop{\mathcal{H}\!\mathit{om}}\nolimits}
\def\SheafExt{\mathop{\mathcal{E}\!\mathit{xt}}\nolimits}
\def\Sch{\mathit{Sch}}
\def\Mor{\mathop{\mathrm{Mor}}\nolimits}
\def\Ob{\mathop{\mathrm{Ob}}\nolimits}
\def\Sh{\mathop{\mathit{Sh}}\nolimits}
\def\NL{\mathop{N\!L}\nolimits}
\def\CH{\mathop{\mathrm{CH}}\nolimits}
\def\proetale{{pro\text{-}\acute{e}tale}}
\def\etale{{\acute{e}tale}}
\def\QCoh{\mathit{QCoh}}
\def\Ker{\mathop{\mathrm{Ker}}}
\def\Im{\mathop{\mathrm{Im}}}
\def\Coker{\mathop{\mathrm{Coker}}}
\def\Coim{\mathop{\mathrm{Coim}}}

% Boxtimes
%
\DeclareMathSymbol{\boxtimes}{\mathbin}{AMSa}{"02}

%
% Macros for moduli stacks/spaces
%
\def\QCohstack{\mathcal{QC}\!\mathit{oh}}
\def\Cohstack{\mathcal{C}\!\mathit{oh}}
\def\Spacesstack{\mathcal{S}\!\mathit{paces}}
\def\Quotfunctor{\mathrm{Quot}}
\def\Hilbfunctor{\mathrm{Hilb}}
\def\Curvesstack{\mathcal{C}\!\mathit{urves}}
\def\Polarizedstack{\mathcal{P}\!\mathit{olarized}}
\def\Complexesstack{\mathcal{C}\!\mathit{omplexes}}
% \Pic is the operator that assigns to X its picard group, usage \Pic(X)
% \Picardstack_{X/B} denotes the Picard stack of X over B
% \Picardfunctor_{X/B} denotes the Picard functor of X over B
\def\Pic{\mathop{\mathrm{Pic}}\nolimits}
\def\Picardstack{\mathcal{P}\!\mathit{ic}}
\def\Picardfunctor{\mathrm{Pic}}
\def\Deformationcategory{\mathcal{D}\!\mathit{ef}}
\usepackage{cleveref}
\crefname{originallemma}{original lemma}{original lemmas}
\Crefname{originallemma}{Original lemma}{Original lemmas}
\makeatletter
\newlabel{morphisms-lemma-locally-finite-type-characterize}{{01T2}{source}{}{sourcepointer-1}{}}
\newlabel{morphisms-lemma-locally-finite-type-characterize@cref}{{[originallemma][0][]01T2}{[1][0][]source}{}{}{}}
\makeatother
\crefname{originallemma}{original lemma}{original lemmas}
\Crefname{originallemma}{Original lemma}{Original lemmas}
\makeatletter
\newlabel{properties-lemma-finite-type-module}{{01PB}{source}{}{sourcepointer-2}{}}
\newlabel{properties-lemma-finite-type-module@cref}{{[originallemma][0][]01PB}{[1][0][]source}{}{}{}}
\makeatother
\crefname{originallemma}{original lemma}{original lemmas}
\Crefname{originallemma}{Original lemma}{Original lemmas}
\makeatletter
\newlabel{algebra-lemma-support-closed}{{00L2}{source}{}{sourcepointer-3}{}}
\newlabel{algebra-lemma-support-closed@cref}{{[originallemma][0][]00L2}{[1][0][]source}{}{}{}}
\makeatother
\crefname{originallemma}{original lemma}{original lemmas}
\Crefname{originallemma}{Original lemma}{Original lemmas}
\makeatletter
\newlabel{morphisms-lemma-support-finite-type}{{056J}{source}{}{sourcepointer-4}{}}
\newlabel{morphisms-lemma-support-finite-type@cref}{{[originallemma][0][]056J}{[1][0][]source}{}{}{}}
\makeatother
\title{A finite type module over an affine finite type scheme is pushed forward from its support by a finitely presented closed immersion}
\author{The Stacks Project Authors}
\date{stacks-project:flat.tex:local-snapshot-20261011; GNU FDL 1.2 or later}
\begin{document}
\maketitle
\noindent\textit{Editorial note.} All delivered mathematical text is the author's own, copied from the pinned Stacks Project chapter file flat.tex; no statement or proof step was rewritten, reordered or removed.\par\medskip
\noindent\textit{Editorial note.} The preamble is the chapter's own retained minimal preamble (sources/stacks-flat/retained-preamble.tex), a verbatim excerpt of the native Xy-pic setup and of the native theorem declarations and mathematical macros; no macro was invented or redefined.\par\medskip
\noindent\textit{Editorial note.} All four referenced results live in earlier chapters, so each locator points at the verbatim line of the project's own tag list staged in the snapshot (sources/stacks-chow/tags.txt) that carries the canonical tag.\par\medskip
\noindent\textit{Allowed standard prerequisites (editorial).}\begin{itemize}
\item Support of a finite type module is the closed subscheme cut out by its annihilator. Algebra, Lemma labelled algebra-lemma-support-closed (canonical Stacks tag 00L2); invoked at proof line 468. Supplies the description of Supp(F) via the annihilator ideal I of M; the comparison of this description with the fibre is the counted step.
\item Characterisation of a finite type morphism of affine schemes as a finite type ring map. Morphisms, Lemma morphisms-lemma-locally-finite-type-characterize (tag 01T2, line 461) and Properties, Lemma properties-lemma-finite-type-module (tag 01PB, line 463). Supply that A -> B is a finite type ring map and that M is a finite B-module, which is what allows finitely many annihilator elements to be chosen later.
\item Support of a finite type module on a fibre. Morphisms, Lemma labelled morphisms-lemma-support-finite-type (canonical Stacks tag 056J); invoked at proof line 474. Supplies that Supp(F\_s) = V(I(B tensor\_A kappa(p))); the choice of f\_1, ..., f\_m cutting out that closed set is inside the counted span.
\end{itemize}\medskip
\ifdefined\hypertarget\hypertarget{sourcepointer-1}{}\fi
\noindent\textit{Source reference.} Original lemma 01T2: sources/stacks-chow/tags.txt, source lines 2359--2359. Cross-chapter standard prerequisite invoked once at proof line 461 to turn the finite type morphism of affines into a finite type ring map; it is a general equivalence and replaces nothing in the retained argument.\par\medskip
\ifdefined\hypertarget\hypertarget{sourcepointer-2}{}\fi
\noindent\textit{Source reference.} Original lemma 01PB: sources/stacks-chow/tags.txt, source lines 2224--2224. Cross-chapter standard prerequisite invoked once at proof line 463 to know that M is a finite B-module; the choice of the annihilator generators is the counted step and stays in the delivered body.\par\medskip
\ifdefined\hypertarget\hypertarget{sourcepointer-3}{}\fi
\noindent\textit{Source reference.} Original lemma 00L2: sources/stacks-chow/tags.txt, source lines 772--772. Cross-chapter standard prerequisite invoked once at proof line 468 for the annihilator description of the support of a finite module.\par\medskip
\ifdefined\hypertarget\hypertarget{sourcepointer-4}{}\fi
\noindent\textit{Source reference.} Original lemma 056J: sources/stacks-chow/tags.txt, source lines 6712--6712. Cross-chapter standard prerequisite invoked once at proof line 474 for the description of the support of the restriction of a finite type module to a fibre.\par\medskip
\begin{lemma}
\label{lemma-sheaf-lives-on-subscheme}
Let $f : X \to S$ be a finite type morphism of affine schemes.
Let $\mathcal{F}$ be a finite type quasi-coherent $\mathcal{O}_X$-module.
Let $x \in X$ with image $s = f(x)$ in $S$.
Set $\mathcal{F}_s = \mathcal{F}|_{X_s}$.
Then there exist a closed immersion $i : Z \to X$ of finite presentation,
and a quasi-coherent finite type $\mathcal{O}_Z$-module $\mathcal{G}$
such that $i_*\mathcal{G} = \mathcal{F}$ and
$Z_s = \text{Supp}(\mathcal{F}_s)$.
\end{lemma}

\begin{proof}
Say the morphism $f : X \to S$ is given by the ring map
$A \to B$ and that $\mathcal{F}$ is the quasi-coherent sheaf
associated to the $B$-module $M$. By
Morphisms, Lemma \ref{morphisms-lemma-locally-finite-type-characterize}
we know that $A \to B$ is a finite type ring map, and by
Properties, Lemma \ref{properties-lemma-finite-type-module}
we know that $M$ is a finite $B$-module. In particular the
support of $\mathcal{F}$ is the closed subscheme of $\Spec(B)$
cut out by the annihilator
$I = \{x \in B \mid xm = 0\ \forall m \in M\}$ of $M$, see
Algebra, Lemma \ref{algebra-lemma-support-closed}.
Let $\mathfrak q \subset B$ be the prime ideal corresponding to $x$
and let $\mathfrak p \subset A$ be the prime ideal corresponding to $s$.
Note that $X_s = \Spec(B \otimes_A \kappa(\mathfrak p))$ and
that $\mathcal{F}_s$ is the quasi-coherent sheaf associated to the
$B \otimes_A \kappa(\mathfrak p)$ module $M \otimes_A \kappa(\mathfrak p)$. By
Morphisms, Lemma \ref{morphisms-lemma-support-finite-type}
the support of $\mathcal{F}_s$ is equal to
$V(I(B \otimes_A \kappa(\mathfrak p)))$. Since
$B \otimes_A \kappa(\mathfrak p)$ is of finite type over $\kappa(\mathfrak p)$
there exist finitely many elements $f_1, \ldots, f_m \in I$
such that
$$
I(B \otimes_A \kappa(\mathfrak p)) =
(f_1, \ldots, f_n)(B \otimes_A \kappa(\mathfrak p)).
$$
Denote $i : Z \to X$ the closed subscheme cut out by
$(f_1, \ldots, f_m)$, in a formula $Z = \Spec(B/(f_1, \ldots, f_m))$.
Since $M$ is annihilated by $I$ we can think of $M$ as an
$B/(f_1, \ldots, f_m)$-module. In other words, $\mathcal{F}$ is the
pushforward of a finite type module on $Z$.
As $Z_s = \text{Supp}(\mathcal{F}_s)$ by construction, this
proves the lemma.
\end{proof}

\end{document}
